% Section 1. Budget 1.75 pages including title and abstract.
% Source: draft/bie-enclosure/s-introduction.tex (clean version).
\section{Introduction}\label{sec:intro}
A line defect in a two-dimensional photonic crystal breaks the periodicity
across the defect but keeps it along the defect. At frequencies in a band gap
of the surrounding crystal, the defect may support waves that propagate along
it and decay exponentially across it, the \emph{guided modes}
\cite{JoannopoulosJohnsonWinnMeade2008}. For a fixed Bloch wavenumber $\beta$
along the defect they are eigenvalues of a self-adjoint operator $A(\beta)$ on
an unbounded strip, located in a gap of its essential spectrum
\cite{Fliss2013}.

Two difficulties arise. The first is that the strip is unbounded across the
defect. Supercell methods truncate it with periodic conditions and incur a
model error that grows as the mode becomes less localized
\cite{Fliss2013,KlindworthSchmidtFliss2014}. Joly, Li, and Fliss constructed
exact boundary conditions for periodic half guides from cell problems
\cite{JolyLiFliss2006}, and Fliss used them to reduce the guided-mode problem to
a bounded neighborhood of the defect, at the price of a nonlinear eigenvalue
problem \cite{Fliss2013}; high-order finite element versions, Robin-to-Robin
variants, and dispersion-curve tracing followed
\cite{KlindworthSchmidtFliss2014,FlissKlindworthSchmidt2015,Klindworth2015}, and
related cell-DtN reductions are used in \cite{HuangLuLi2007,YuanLu2007}. The
second difficulty is reliability. For piecewise constant media, boundary
integral equations discretize only the interfaces and, with quasi-periodic Green
functions, reach high order with few unknowns
\cite{Linton1998,BarnettGreengard2010,HaoBarnettMartinssonYoung2014}. A nearly
singular matrix in a nonlinear eigenvalue problem, however, only shows that a
finite-dimensional equation is nearly solvable \cite{GuettelTisseur2017}; it
neither bounds the distance to the spectrum of the continuous operator nor
proves that a spectral point exists.

\emph{Related enclosure methods.} Bounding eigenvalues by the boundary defect of
trial fields that solve the equation exactly is the classical method of
particular solutions \cite{MolerPayne1968,BetckeTrefethen2005}. Descloux and
Tolley extended it to domain decomposition, with local exact solutions and an
eigenvalue error controlled by the mismatch of values and gradients between
subdomains \cite{DesclouxTolley1983} (see also \cite[Lem.~2.1]{Betcke2007});
this is the closest structure to ours. Barnett, Hassell, and Tacy bounded the
distance to the spectrum by boundary residuals for Neumann and Dirichlet
problems \cite{BarnettHassell2011,BarnettHassellTacy2018}. Dahne and Salvy
computed candidates by any means and then bounded the residual and the interior
norm with interval arithmetic and Taylor models \cite{DahneSalvy2020}; we
separate search and reliable evaluation in the same way. Temple, Kato--Temple,
and Lehmann--Goerisch bounds \cite{Temple1928,NakaoPlumWatanabe2019} require
trial functions in the domain of the operator or its form, which our
piecewise field is not. These methods do handle essential spectrum
\cite{ZimmermannMertins1995,DaviesPlum2004}, infinite waveguides
\cite{BehnkeMertinsPlumWieners2000}, and photonic-crystal line defects
\cite{BrownHoangPlumWood2009,HoangPlumWieners2009}; the finite element comparison
in \cref{sec:ex1} is of this type. All the particular-solution results above
concern bounded domains. Here the strip is unbounded, the discrete field is
generated cell by cell in the half guides, and besides the material interfaces
there are infinitely many artificial walls and seams.

\emph{Contributions.} We consider a strip with periodic media on both sides,
possibly different, and a local defect in a center cell, with piecewise
constant coefficients and one rod per cell. The main result is a
spectral-distance theorem for piecewise Helmholtz fields on this unbounded
strip (\cref{thm:distance}): for a trial field that solves the Helmholtz
equation exactly in every homogeneous piece,
$\operatorname{dist}(\widehat k^2,\sigma(A(\beta)))$ is bounded by the grouped
$L^2$ norms $\varepsilon_g$ of its jumps across rod boundaries and across the
infinitely many artificial walls and seams, weighted by trace constants $C_g$
that are explicit in the geometry. The resolvent mechanism behind such bounds
is classical; the theorem adapts it to a field generated cell by cell in both
half guides, with constants that can be evaluated. The second result is a
finite-dimensional test for a common gap of the two leads
(\cref{thm:complement,prop:panel}). A comparison with the Fourier complement
gives guaranteed lower bounds for the band above the gap from a small
plane-wave basis, and a Lipschitz bound in the Bloch wavenumber turns finitely
many inertia tests into bounds for all wavenumbers. Third, every input is
evaluated with rounding included: the jump integrals on curved interfaces, the
sums over infinitely many lead cells through the propagation matrix, the mass
lower bound, and the gap test. Together these give an interval that contains an
eigenvalue by the common gap alone, without a local isolation condition or a
mode count. For a star-shaped defect between two different leads, one run of
$14$ minutes gives a $\mu$-interval of width $2.74\times10^{-11}$; for modes
approaching a gap edge, where the lead contraction rate rises to $0.80$, the
width stays below $10^{-11}$.

Candidates are computed by boundary integral equations assembled from known
components: transparent boundary conditions for the half guides
\cite{JolyLiFliss2006,Fliss2013,FlissKlindworthSchmidt2015}, Rayleigh
expansions at artificial interfaces, and near-image splitting and proxy
representations of quasi-periodic fields
\cite{BarnettGreengard2010,BarnettGreengard2011,ChoBarnett2015}. None of these
is claimed as new, and the enclosure does not depend on the accuracy of the
candidate. We do not claim that the eigenvalue in the interval is unique, nor
its index or multiplicity, and we do not study a priori convergence: the
theorem applies to a given candidate and field, however they were obtained.

\Cref{sec:model} introduces the model and the gap test, \cref{sec:discretization}
the computation of candidates, and \cref{sec:bound} the trial field, the
spectral-distance theorem and the rounding-inclusive evaluation.
\Cref{sec:results} reports the examples, including a comparison with a finite
element enclosure, and \cref{sec:conclusion} concludes. The results stated here
are proved in the main text; the supplementary materials contain the error
bounds of the reliable evaluation and additional numerical data.
